\section{Introduction}

\subsection{A multiple flag variety and its orbital decomposition}
\label{section-1.1}
In this paper, we consider the multiple flag variety
\begin{equation}
\label{equation-dfv}
\Xfv=\Grass(V,r)\times \Flags(V^+)\times \Flags(V^-),
\end{equation}
where
\begin{itemize}
\item $V=\mC^{p+q}$ is equipped with a polar decomposition $V=V^+\oplus V^-$ with $V^+=\mC^p\times\{0\}^q$ and $V^-=\{0\}^p\times \mC^q$;
\item $\Grass(V,r)$ denotes the Grassmann variety of $r$-dimensional subspaces of $V$;
\item $\Flags(V^+)$ and $\Flags(V^-)$ denote the varieties of complete flags of $V^+$ and $V^-$, respectively.
\end{itemize}
Each factor of the variety $\Xfv$ has a natural action of
\[
K\coloneqq \GL(V^+)\times \GL(V^-)=\left\{\begin{pmatrix} a & 0\\ 0 & d \end{pmatrix}:a\in\GL_p(\mC),\ d\in\GL_q(\mC)\right\}\subset\GL(V),
\]
and $\Xfv$ is endowed with the resulting diagonal action of $K$.

The multiple flag variety $\Xfv$ can be written in the form
\[
\Xfv=G/P\times K/B_K,
\]
where $G=\GL(V)$, $P\subset G$ is a maximal parabolic subgroup, $B_K\subset K$ is a Borel subgroup.
In this way, $\Xfv$ is a \textit{double flag variety associated to the symmetric pair}
$(G,K)$
in the sense of \cite{NO} and \cite{HNOO}.
In particular, it is known from \cite{NO} that the above variety $\Xfv$ has a finite number of $K$-orbits.

In \cite{FN1,FN2}, we have initiated an analogue of Steinberg theory for double flag varieties associated to symmetric pairs such as $\Xfv$.
Specifically, we have defined two Steinberg maps, from the set of $K$-orbits of $\Xfv$ to the sets of nilpotent $K$-orbits of $\fk\coloneqq \Lie(K)$ and of its Cartan complement $\mathfrak{s}$, respectively.

For general symmetric pairs, the calculation of the generalized Steinberg maps appears to be quite difficult.
In \cite{FN2}, we have considered the variety $\Xfv$ of \eqref{equation-dfv} in the special case where $p=q=r$ and we have computed the Steinberg maps on a special subset of $K$-orbits parametrized by partial permutations.
The present paper deals with the variety $\Xfv$ of \eqref{equation-dfv} and its $K$-orbits in full generality.

Here we summarize the main results achieved in this paper:
\begin{itemize}
\item We describe completely the decomposition of $\Xfv$ into $K$-orbits: we give a para\-metrization of the orbits (in terms of certain graphs), we provide a dimension formula, and describe the closure relations and the cover relations; see Section \ref{section-2.2}.
\item Our main result (Theorem \ref{T2}) is the calculation of the two Steinberg maps mentioned above. This calculation is concrete, and it is done by means of a sophisticated combinatorial algorithm that generalizes the Robinson--Schensted correspondence; see Sections \ref{section-2.3}--\ref{section-2.4}.
\end{itemize}
In the following subsection, we explain the construction of the generalized Steinberg maps and we give more insight on our main result.

\subsection{Conormal variety and Steinberg maps}
\label{section-1.2}
We consider the Lie algebras
\[
\fg=\gl_{p+q}(\mC)\supset\fk=\left\{\begin{pmatrix} a & 0\\ 0 & d \end{pmatrix}:a\in\gl_p(\mC),\ d\in\gl_q(\mC)\right\}
\]
and a Cartan decomposition
\[
\fg=\fk\oplus\fs\qquad\mbox{where}\qquad\fs=\left\{\begin{pmatrix} 0 & b\\ c & 0 \end{pmatrix}:b\in\Mat_{p,q}(\mC),\ c\in\Mat_{q,p}(\mC)\right\}.
\]
We write $x=x_\fk+x_\fs$ with $(x_\fk,x_\fs)\in\fk\times\fs$ for the decomposition of an element $x\in\fg$ along the Cartan decomposition.
Moreover, we identify the Lie algebras $\fg$ and $\fk$ with their duals $\fg^*$ and $\fk^*$ through the trace form.

Any partial flag $\flag=(F_0=0\subset F_1\subset\ldots\subset F_k=V')$ of a vector space $V'$ gives rise to a parabolic subalgebra $\fp(\flag)$ of $\gl(V')$ and the corresponding nilradical $\nil(\flag)$ defined by
\begin{eqnarray*}
\fp(\flag) & = & \mathrm{Stab}_{\gl(V')}(\flag)=\{x\in\gl(V'):x(F_i)\subset F_i\quad \forall i=1,\ldots,k\},\\
\nil(\flag) & = & \{x\in\gl(V'):x(F_i)\subset F_{i-1}\quad\forall i=1,\ldots,k\}.
\end{eqnarray*}
For a subspace $W\subset V'$, we denote by $\fp(W)$ and $\nil(W)$ the (maximal) parabolic subalgebra and the nilradical associated to the partial flag $(0\subset W\subset V')$.

As explained in \cite[\S3]{FN2}, the cotangent bundle $T^*\Xfv$ inherits a Hamiltonian action of $K$, which gives rise to a moment map $\mu_\Xfv:T^*\Xfv\to\fk^*=\fk$.
The nullfiber $\conormal=\mu_\Xfv^{-1}(0)$ is called a {\em conormal variety}. It can be described explicitly as
\[
\conormal=\{(W,\flag^+,\flag^-,x)\in \Xfv\times\gl(V):x\in\nil(W),\ x_\fk\in\nil(\flag^+)\times\nil(\flag^-)\}.
\]

Every $K$-orbit $\Xorbit\subset \Xfv$ yields a conormal bundle $T^*_\Xorbit\Xfv$ that can be realized as a (locally-closed) subvariety of $\conormal$ given by
\[
T^*_\Xorbit\Xfv=\{(W,\flag^+,\flag^-,x)\in\conormal:(W,\flag^+,\flag^-)\in\Xorbit\}.
\]
The variety $\conormal$ is equidimensional of dimension $\dim\Xfv$, and its irreducible components are precisely the closures of the various conormal bundles $T^*_\Xorbit\Xfv$, since the set of orbits $\Xfv/K$ is finite.
One can find a more comprehensive introduction to the theory of conormal varieties in \cite[\S 3]{FN2} (see also \cite{CG}).

The conormal variety $\conormal$ is equipped with two $K$-equivariant projections to $\fk$ and $\fs$, namely
\[\phi_\fk:\conormal\to\fk,\ (W,\flag^+,\flag^-,x)\mapsto x_\fk\quad\mbox{and}\quad
\phi_\fs:\conormal\to\fs,\ (W,\flag^+,\flag^-,x)\mapsto x_\fs.\]
It immediately follows from the description of the conormal variety that the image of $\phi_\fk$ is contained in the cone of nilpotent elements $\nilpotentsof{\fk}\subset\fk$.
It is shown in \cite[Proposition 4.2]{FN1} that (for the variety $\Xfv$ considered in this paper) the image of $\phi_\fs$ is also contained in the nilpotent cone $\nilpotentsof{\fs}\subset\fs$.
It is known that both nilpotent cones $\nilpotentsof{\fk}$ and $\nilpotentsof{\fs}$ consist of finitely many adjoint $K$-orbits

Therefore, we can define two maps
\[\Phi_\fk:\Xfv/K\to\nilpotentsof{\fk}/K
\quad\mbox{and}\quad
\Phi_\fs:\Xfv/K\to\nilpotentsof{\fs}/K\]
in the following way: for every orbit $\Xorbit\in\Xfv/K$, define $\Phi_\fk(\Xorbit)\in \nilpotentsof{\fk}/K$, resp. $\Phi_\fs(\Xorbit)\in\nilpotentsof{\fs}/K$, as the unique nilpotent $K$-orbit which is open and dense in the image of the conormal bundle $T^*_\Xorbit\Xfv$ by the projection map $\phi_\fk$, resp. $\phi_\fs$. According to the terminology introduced in \cite{FN2}, we will refer to $\Phi_\fk$ as the {\em symmetrized Steinberg map} and to $\Phi_\fs$ as the {\em exotic Steinberg map}.

In Section \ref{section-2.3}, we describe the maps $\Phi_\fk$ and $\Phi_\fs$. In \cite{FN2}, in the special case where $p=q=r$,
the images $\Phi_\fk(\Xorbit)$ and $\Phi_\fs(\Xorbit)$ are determined when an orbit $\Xorbit$ is contained in a ``big cell'' of $\Xfv/K$.
Moreover, in Section \ref{section-2.4}, we describe the fibers of $\Phi_\fk$ by means of a combinatorial procedure that extends the classical Robinson--Schensted correspondence; this also generalizes \cite[Theorem 7.8]{FN2}.
In this way, the results in the present paper are new and complement those in \cite{FN2}, as now we have a full description of the two Steinberg maps and a better understanding of them at the same time. Note that the results given in Section \ref{section-2.3} regarding $\Phi_\fk$ and for $p=q=r$ were already announced in \cite[\S 2]{FN3} without proofs.




\section{Main results}

\subsection{Combinatorial notation on pairs of partial permutations}

\label{section-2.1}

By $\ppermutationsof_{p,r}$ we denote the set of $p\times r$ matrices whose coefficients are $0$ or $1$, with at most one $1$ in each row and each column.
(If $p=r$, we recover the set of partial permutation matrices considered in \cite{FN2}.)
By $\ppermutationsof=\ppermutationsof_{(p,q),r}$ we denote the set of $(p+q)\times r$ matrices of rank $r$ (we have $r\leq p+q$) of the form
\[\omega=\begin{pmatrix} \tau_1\\ \tau_2 \end{pmatrix},\]
where $\tau_1\in\ppermutationsof_{p,r}$ and $\tau_2\in\ppermutationsof_{q,r}$.
Note that the symmetric group $\permutationsof{r}$ acts on $\ppermutationsof$ by right multiplication, and we denote the quotient set by $\parameters=\ppermutationsof/\permutationsof{r}$.

{\em In Section \ref{section-2.2}, we will show that the elements of $\parameters$ parameterize the $K$-orbits of~$\Xfv$.}


\paragraph{\bf a) Graphic representation of a pair of partial permutations:}\

We represent any element $\omega\in\parameters$ by a graph $\mathcal{G}(\omega)$ obtained as follows:
\begin{itemize}
\item
The set of vertices consists of $p$ ``positive'' vertices $1^+,\ldots,p^+$ and $q$ ``negative'' vertices $1^-,\ldots,q^-$, displayed along two horizontal lines.
%The set of vertices $\mathcal{V}^+\cup\mathcal{V}^-$ consists of ``positive'' vertices $\mathcal{V}^+=\{1^+,\ldots,p^+\}$ and ``negative'' vertices $\mathcal{V}^-=\{1^-,\ldots,q^-\}$, displayed along two horizontal lines.
\item Put an edge between $i^+$ and $j^-$ for every column of $\omega$ that contains exactly two $1$'s, in positions $i$ (within the block $\tau_1$) and $p+j$ (within the block $\tau_2$).
\item
Put a mark at the vertex $i^+$, respectively $j^-$, for every column of $\omega$ that contains exactly one $1$, in position $i$ (within $\tau_1$), respectively $p+j$ (within $\tau_2$).
\end{itemize}
For instance, for $(p,q)=(5,3)$ and $r=4$,
\begin{equation}
\label{2.1}
\omega=\mbox{\tiny $\begin{pmatrix} 0 & 0 & 0 & 0\\ 1 & 0 & 0 & 0\\ 0 & 0 & 0 & 0\\ 0 & 1 & 0 & 0\\ 0 & 0 & 1 & 0\\ \hline 0 & 1 & 0 & 0\\ 0 & 0 & 0 & 1\\ 1 & 0 & 0 & 0 \end{pmatrix}$}
\qquad
\Rightarrow
\qquad
\mathcal{G}(\omega)=
\mbox{\tiny
$\begin{picture}(100,20)(0,0)
\put(0,11){$\bullet$}\put(20,11){$\bullet$}\put(40,11){$\bullet$}\put(60,11){$\bullet$}\put(80,11){$\bullet$}
\put(0,18){$1^+$}\put(20,18){$2^+$}\put(40,18){$3^+$}\put(60,18){$4^+$}\put(80,18){$5^+$}
\put(0,-11){$\bullet$}\put(20,-11){$\bullet$}\put(40,-11){$\bullet$}
\put(0,-20){$1^-$}\put(20,-20){$2^-$}\put(40,-20){$3^-$}
\put(21,13){\line(1,-1){21}}
\put(61,12){\line(-3,-1){60}}
\put(82,13){\circle{8}}
\put(22.5,-9){\circle{8}}
\end{picture}$}.
\end{equation}
We will say that a vertex is {\em free} whenever it is not a marked
point nor an end point of an edge (like $1^+$ and $3^+$ in the above
example).

In general, the assignment $\omega\mapsto\mathcal{G}(\omega)$ establishes a bijection between $\parameters$ and the set of graphs with vertices %$\mathcal{V}^+\cup\mathcal{V}^-$,
$\{1^+,\ldots,p^+\}\cup\{1^-,\ldots,q^-\}$,
exactly $r$ edges or marked vertices, where every vertex is incident with at most one edge, and such that there is no edge which is incident with a marked vertex or joins two vertices of the same sign.


\paragraph{\bf b) Numerical invariants for the graph:}\

The following data associated to an element $\omega\in\parameters$ and its graphic representation $\mathcal{G}(\omega)$ will play a role in the statement of our main results.
\begin{itemize}
\item Set the {\em degree} of a vertex of $\mathcal{G}(\omega)$ as $0$, $1$, or $2$, depending on whether this vertex is free, incident with an edge, or marked.

We define $a^+(\omega)$, respectively $a^-(\omega)$, as the number of pairs
of positive vertices $(i^+,j^+)$ with $i<j$ and $\deg(i^+)<\deg(j^+)$,
respectively pairs of negative vertices $(i^-,j^-)$ with $i<j$ and $\deg(i^-)<\deg(j^-)$.

Let $b(\omega)$ be the number of edges of $\mathcal{G}(\omega)$.

Finally, let $c(\omega)$ be the number of \textit{crossings}, i.e. pairs of edges $(i^+,j^-)$, $(k^+,\ell^-)$ such that $i<k$ and $j>\ell$.

\item For all $(i,j)\in\{0,1,\ldots,p\}\times\{0,1,\ldots,q\}$, let $r_{i,j}(\omega)$ be the number of edges and marks contained in the subgraph of $\mathcal{G}(\omega)$ formed by the vertices $k^+$ ($1\leq k\leq i$), $\ell^-$ ($1\leq \ell\leq j$) and the edges/marks contained within this set of vertices. Let $R(\omega)=\left( r_{i,j}(\omega) \right)_{0\leq i\leq p,\,0\leq j\leq q}$ be the $(p+1)\times(q+1)$ matrix containing these numbers.

In particular, $r_{i,0}(\omega)$ (resp. $r_{0,j}(\omega)$) is the
number of marked vertices among
$\{1^+,\ldots,i^+\}$ (resp. $\{1^-,\ldots,j^-\}$).
\end{itemize}
{\em In Section \ref{section-2.2}, the numbers $a^\pm(\omega)$, $b(\omega)$, $c(\omega)$ appear in the dimension formula for the $K$-orbits of $\Xfv$, while the matrices $R(\omega)$ are used to describe the inclusion relations between orbit closures.}
\begin{itemize}
\item We decompose $\{1,\ldots,p\}=I\sqcup L\sqcup L'$ in the following way:
$I$, resp. $L$, resp. $L'$, denotes the set of elements $i\in\{1,\ldots,p\}$ such that $i^+$ is a vertex of $\mathcal{G}(\omega)$ of degree $1$, resp. $2$, resp. $0$.

We decompose $\{1,\ldots,q\}=J\sqcup M\sqcup M'$ in the same way: $J$, resp. $M$, resp. $M'$, consists of the elements $j$ such that $j^-$ has degree $1$, resp. $2$, resp. $0$.

Let $\sigma:J\to I$ be the bijection defined by letting $\sigma(j)=i$ if $(i^+,j^-)$ is an edge in $\mathcal{G}(\omega)$.
\end{itemize}
Note that $\omega$ is characterized by the subsets $I$, $L$, $L'$, $J$, $M$, $M'$ and the bijection $\sigma:J\to I$.
{\em In Section \ref{section-2.3},
these data are used to compute
the symmetrized and exotic Steinberg maps, by means of a combinatorial algorithm.}

Note also that we have $b(\omega)=\card{I}=\card{J}$, and $c(\omega)$ is the number of inversions of~$\sigma$.



\begin{exam}
\label{E2.1}
Let $\omega$ be as in \eqref{2.1}. Then,
\begin{eqnarray*}
 & a^+(\omega)=7,\quad a^-(\omega)=1,\quad b(\omega)=2,\quad c(\omega)=1, \quad
R(\omega)=\mbox{\tiny $\begin{pmatrix} 0 & 0 & 1 & 1\\ 0 & 0 & 1 & 1\\ 0 & 0 & 1 & 2\\ 0 & 0 & 1 & 2\\ 0 & 1 & 2 & 3\\ 1 & 2 & 3 & 4 \end{pmatrix}$}\\
&
\begin{array}{llll}
I=\{2,4\}, & L=\{5\}, & L'=\{1,3\},\\[2mm]
J=\{1,3\}, & M=\{2\}, & M'=\emptyset,
\quad\sigma=\begin{pmatrix} 1 & 3\\ 4 & 2 \end{pmatrix}\in\mathrm{Bij}(J,I).
\end{array}
\end{eqnarray*}
Note that the matrix $R(\omega)$ can also be viewed as a plane partition.
\end{exam}

\subsection{Orbit decomposition of the multiple flag variety $\Xfv$}

\label{section-2.2}

Recall that we consider the space $V=\mC^{p+q}$ endowed with the polar decomposition
\[V=V^+\oplus V^-\qquad\mbox{where}\qquad V^+=\mC^p\times\{0\}^q\quad\mbox{and}\quad V^-=\{0\}^p\times \mC^q.\]
Let
\[
\flag_0^+=(\mC^i\times\{0\}^{p-i}\times\{0\}^q)_{i=0}^p
\quad\mbox{and}\quad
\flag_0^-=(\{0\}^{p}\times\mC^j\times\{0\}^{q-j})_{j=0}^q
\]
be the standard complete flags of $V^+$ and $V^-$.

Every  $(p+q)\times r$ matrix $\omega$ determines a subspace $[\omega]\coloneqq \mathrm{Im}\,\omega\subset V$, which remains the same up to permutation of the columns of $\omega$.
In particular, every $\omega\in\parameters$ determines a point $[\omega]$ in $\Grass(V,r)$, and thus a point $([\omega],\flag_0^+,\flag_0^-)$ in $\Xfv=\Grass(V,r)\times\Flags(V^+)\times\Flags(V^-)$.

\begin{theorem}
\begin{enumerate}[label=(\arabic*)]
\item\label{Tp11} Every $K$-orbit in $\Xfv$ is of the form $\Xorbit_\omega\coloneqq K\cdot([\omega],\flag_0^+,\flag_0^-)$ for a unique element $\omega\in\parameters=\ppermutationsof_{(p,q),r}/\permutationsof{r}$.
\item \label{Tp12} $\dim\Xorbit_\omega=\frac{p(p-1)}{2}+\frac{q(q-1)}{2}+a^+(\omega)+a^-(\omega)+\frac{b(\omega)(b(\omega)+1)}{2}+c(\omega)$.
\item \label{Tp13} $\Xorbit_\omega$ is the set of triples $(W,\flag^+=(F^+_i)_{i=0}^p,\flag^-=(F_j^-)_{j=0}^q)\in\Xfv$
satisfying the condition
\[\dim W\cap(F_i^++F_j^-)=r_{i,j}(\omega)\ \ \mbox{for all $(i,j)\in\{0,\ldots,p\}\times\{0,\ldots,q\}$}.\]
 \item \label{Tp14} $\overline{\Xorbit_\omega}\subset\overline{\Xorbit_{\omega'}}$ if and only if $r_{i,j}(\omega)\geq r_{i,j}(\omega')$ for all $(i,j)\in\{0,\ldots,p\}\times\{0,\ldots,q\}$.
\end{enumerate}
\label{Tp1}
\end{theorem}

As a complement of this result, we determine the cover relations in the poset $(\{\overline{\Xorbit_\omega}\},\subset)$. We say that $\Xorbit_{\omega'}$ covers $\Xorbit_\omega$ if $\overline{\Xorbit_{\omega'}}$ strictly contains $\overline{\Xorbit_\omega}$ and is minimal (among the orbit closures) for this property.
Equivalently, this means that $\overline{\Xorbit_\omega}$ is an irreducible component of the boundary
$\partial\Xorbit_{\omega'}=\overline{\Xorbit_{\omega'}}\setminus \Xorbit_{\omega'}$.

\begin{theorem}
\label{C1}
The following conditions are equivalent:
\begin{enumerate}
\item $\Xorbit_{\omega'}$ covers $\Xorbit_\omega$;
\item $\dim\Xorbit_{\omega'}=\dim\Xorbit_\omega+1$ and (the graph of) $\omega$ is obtained from (the graph of) $\omega'$ by modifying the pattern of at most four vertices $a^+,b^+,c^-,d^-$ ($a<b$, $c<d$), according to one of the cases indicated in Figure \ref{figure1}.
\end{enumerate}

As a consequence, the boundary of every non-closed orbit is equidimensional of codimension one.
\renewcommand{\arraystretch}{1.2}
\begin{figure}[h]
\begin{tabular}{||>{\centering}p{4.2cm}||>{\centering\arraybackslash}p{4.2cm}||}


\hline {\rm Case 1} & {\rm Case 2}\\

\hline

&\\[-2mm]

\coverA$\rightsquigarrow\ $\coverB\ &

\begin{tabular}{c}\coverC $\rightsquigarrow$ \coverD\\[6mm] {\rm or}\\[3mm]

\coverE$\rightsquigarrow\ $\coverF \end{tabular}
\\[20mm] \hline
\end{tabular}

\begin{tabular}{||>{\centering}p{2.98cm}||>{\centering}p{2cm}||>{\centering\arraybackslash}p{2.98cm}||}


\hline {\rm Case 3} & {\rm Case 4} & {\rm Case 5}\\

\hline

& &\\[-2mm]

\begin{tabular}{c}\coverG $\rightsquigarrow$ \coverH\\[6mm] {\rm or}\\[3mm]

\coverI$\rightsquigarrow\ $\coverJ \end{tabular} &

\begin{tabular}{c}\coverK $\rightsquigarrow$ \coverL\\[6mm] {\rm or}\\[3mm]

\coverM$\rightsquigarrow\ $\coverN \end{tabular} &

\begin{tabular}{c}\coverO $\rightsquigarrow$ \coverP\\[6mm] {\rm or}\\[3mm]

\coverQ$\rightsquigarrow\ $\coverR \end{tabular}\\[20mm] \hline

\end{tabular}
\caption{Elementary moves yielding cover relations in the poset $(\{\overline{\Xorbit_\omega}\},\subset)$.}
\label{figure1}
\end{figure}
\renewcommand{\arraystretch}{1}
\end{theorem}

It follows from Theorem \ref{C1} that the boundary $\partial\Xorbit_\omega\coloneqq \overline{\Xorbit_\omega}\setminus \Xorbit_\omega$ of every non-closed orbit is equidimensional of codimension one in $\overline{\Xorbit_\omega}$.
This boundary is in general not irreducible as it already appears in Example \ref{E2.4}\,{\rm (a)}.

\begin{exam}
\label{E2.4}
(a)
In Figure \ref{figure2}, we represent the elements $\omega\in\parameters$ (under the form of their graphic incarnations $\mathcal{G}(\omega)$) in the case where $p=q=r=2$. We indicate the dimensions of the corresponding $K$-orbits $\Xorbit_\omega$. An edge joining two parameters indicates a cover relation.

\begin{figure}[h!]
\textcolor{blue}{\xymatrixcolsep{0.4pc}
\xymatrix{
\textcolor{darkgray}{\dim:6}
& & & & & \graphA \ar@/_/@{-}[lld] \ar@{-}[d] \ar@/^/@{-}[rrd]  & & & & &\\
\textcolor{darkgray}{5} & & & \graphB \ar@{-}[lld] \ar@{-}[d] \ar@{-}[rrd] & & \graphC \ar@/_/@{-}[lllld] \ar@{-}[lld] \ar@{-}[rrd] \ar@/^/@{-}[rrrrd]  & & \graphD \ar@{-}[lld] \ar@{-}[d] \ar@{-}[rrd] & & &\\
\textcolor{darkgray}{4} & \graphE \ar@{-}[rd]\ar@/_/@{-}[rrrd] & & \graphF \ar@{-}[ld]\ar@/_/@{-}[rrrd] & & \graphG \ar@/_/@{-}[ld]\ar@/^/@{-}[rd] & & \graphH \ar@{-}[rd]\ar@/^/@{-}[llld] & & \graphI \ar@{-}[ld]\ar@/^/@{-}[llld] &\\
\textcolor{darkgray}{3} & & \graphJ \ar@{-}[rd]\ar@/_/@{-}[rrrd] & & \graphK \ar@{-}[rd] & & \graphL \ar@{-}[ld] & & \graphM \ar@{-}[ld]\ar@/^/@{-}[llld] & &\\
\textcolor{darkgray}{2} & & & \graphN & & \graphO & & \graphP & & &}}
\caption{The parameters of the $K$-orbits of $\Xfv$ and the cover relations for $p=q=r=2$.}
\label{figure2}
\end{figure}

(b) In Figure \ref{figure1}, the vertices $a^+,b^+$ or $c^-,d^-$ involved in an elementary move that yields a cover relation are not necessarily consecutive. For example, $\Xorbit_{\omega'}$ covers $\Xorbit_\omega$
if
\[
\mathcal{G}(\omega')=\graphnA
\qquad\qquad
\mathcal{G}(\omega)=\graphnB
\]

\bigskip
\noindent
This corresponds to Case 1 of Figure \ref{figure1} with $a=c=1$, $b=d=3$.
Note however that in Case 5 of Figure \ref{figure1}, the vertices $a^+,b^+$ (resp. $c^-,d^-$) must be consecutive for having a cover relation.
\end{exam}

The proofs of Theorems~\ref{Tp1} and \ref{C1} are given in Section \ref{section-3}.

One ingredient for showing the parametrization of the orbits in Theorem \ref{Tp1}\ref{Tp11} is that
the orbits of a pair of Borel subgroups $B_p^+\times B_r^+\subset\GL_p(\mC)\times\GL_r(\mC)$ on the space of $p\times r$ matrices are parametrized by partial permutations
(Lemma \ref{L3.2}).
This classification of orbits is also shown in \cite{Fulton-1991}, where dimension formulas, closure relations, and properties of closures of orbits are described.
